\subsection{This Day You Shall Know: the Vigil}\label{sec:vigil}

On 24 December the station is St Mary Major, and the Mass of the Vigil opens
with \latin{Hodie scietis, quia veniet Dominus, et salvabit nos: et mane
videbitis gloriam ejus}. The words are Moses', in Exodus:

\begin{quote}
In the evening you shall know that the Lord hath brought you forth out of the
land of Egypt: \dots\ in the morning you shall see the glory of the Lord\dots\par
{\smallskip\noindent\small ---Exodus 16:6--7\par}
\end{quote}

Where Moses said \latin{vespere}, in the evening, the Church says \latin{hodie},
this day; and where he spoke of the Lord who had brought them out of Egypt, she
speaks of the Lord who will come and save us.

It is worth reading on in Exodus to see what glory Moses meant. Israel has been
murmuring in the desert, and God has just said to Moses, ``Behold I will rain
bread from heaven for you''. Then:

\begin{quote}
So it came to pass in the evening, that quails coming up, covered the camp: and
in the morning a dew lay round about the camp. And when it had covered the
face of the earth, it appeared in the wilderness small, and as it were beaten
with a pestle, like unto the hoar frost on the ground. And when the children of
Israel saw it, they said one to another: Manhu! which signifieth: What is this!
for they knew not what it was. And Moses said to them: This is the bread which
the Lord hath given you to eat.\par
{\smallskip\noindent\small ---Exodus 16:13--15\par}
\end{quote}

In the morning came the bread, with the dew. The Mass of the Fourth Sunday
began with \latin{Rorate}, ``let the clouds rain the just''; Isaiah's verb, \latin{pluant}, is the verb God uses in
Exodus, \latin{pluam vobis panes de cælo}, ``I will rain bread from heaven''. On
the last day of Advent the Church takes her Introit from Moses' promise of the
manna. The people in the desert looked at the bread and said, ``What is
this!'' Tomorrow the shepherds will look into a manger.

The Gospel of the Vigil is St Matthew's account of Joseph. He finds Mary with
child, and is minded to put her away privately, and an angel comes to him in
his sleep: ``And she shall bring forth a son: and thou shalt call his name Jesus.
For he shall save his people from their sins.'' And the Collect, the prayer of
the Redeemer and the Judge, asks that we who receive him now with joy may see
him without fear when he comes as Judge.

\subsubsection*{The proclamation}

On the Vigil, at Prime, the Breviary reads the Martyrology for the next day,
and the entry for 25 December begins with a proclamation. It counts the ages of
the world to the birth of Christ. In the English of the Breviary it runs:

\begin{quote}
In the year 5199 from the creation of the world, when in the beginning God
created the heavens and the earth; In the year 2959 from the flood; In the year
2015 from the birth of Abraham; In the year 1510 from the going forth of the
people of Israel out of Egypt under Moses; In the year 1032 from the anointing
of David as King; In the 65th week according to the prophecy of Daniel; In the
194 Olympiad; In the 752 from the foundation of the city of Rome; In the 42nd
year of the reign of the Emperor Octavian Augustus; In the 6th age of the world,
while the whole earth was at peace, Jesus Christ, Himself Eternal God and Son of
the Eternal Father, being pleased to hallow the world by His most gracious
coming, having been conceived of the Holy Ghost, and when nine months were
passed after His conception, (all kneel down) was born of the Virgin Mary at
Bethlehem of Juda made Man, our Lord Jesus Christ was born according to the
flesh.\par
{\smallskip\noindent\small ---The Roman Martyrology, 25 December\par}
\end{quote}

Creation, the flood, Abraham, Moses, David, Daniel's weeks, the Olympiads, the
founding of Rome, the reign of Augustus: the ways men had of counting time are
gathered up and brought to one night, and at the words of the birth all
kneel. Like the fifteenth year of Tiberius on the Fourth Sunday, the
proclamation sets the eternal Son inside a calendar.

St Bernard preached on the Vigil, and he began with words very like the
Martyrology's:

\begin{quote}
Jesus Christ, the Son of God, is born in Bethlehem of Juda. What heart so stony
as not to be softened at these words? What soul is not melted at this voice of
her Beloved? What announcement could be sweeter? what intelligence more
enrapturing? Was its like ever heard before? or when did the world ever receive
such tidings? Jesus Christ, the Son of God, is born in Bethlehem of Juda. O
short word, telling of the Eternal Word abbreviated for us! O word full of
heavenly delights!\par
{\smallskip\noindent\small ---St Bernard, First Sermon on the Vigil of the Nativity\par}
\end{quote}

Read the proclamation aloud at home on the morning of the Vigil, all of it,
slowly, and kneel where the Church kneels. Then let the day go toward
Bethlehem with the shepherds: ``Let us go over to Bethlehem and let us see this
word that is come to pass''.
